\documentclass[Typ=PLAIN,Semester=1,Vorlesung=FG]{Lilly}
\usetikzlibrary{decorations.text}

\input Konanur.fd

%% Mathe
\definecolor{FacultyMathexColor}{RGB}{086,170,028}
%% Praktische Informatik
\definecolor{FacultyPraktischeInformatikxColor}{RGB}{221,101,015}
%% Theoretische Informatik
\definecolor{FacultyTheoretischeInformatikxColor}{RGB}{002,133,221}
%% Technische Informatik
\definecolor{FacultyTechnischeInformatikxColor}{RGB}{177,2,221}

\newlength{\LILLYxLeftsep}\setlength{\LILLYxLeftsep}{2cm}

%% Upper Text %% Symbol %% WHITE Color %% BACK Color %% Tikz %% Size %% FONTSIZE UPPER
\DeclareDocumentCommand{\LILLYxGENxFACULTY} { m m m O{} O{} O{} O{} O{0.1} O{0.05}} {
    \begin{tikzpicture}[#5]
        \draw [decorate,decoration={text along path,text align=center,
        text={|\footnotesize\normalNumber\color{#3}#7|#1}}]
       (-1.275-#8,0) arc [start angle=180, end angle=00, radius=1.275+#8];
       \draw(-1.5,0) node {{$\bullet$}} ++ (3,0) node {{$\bullet$}};
       \draw [decorate,decoration={text along path,text align=center,
       text={|\footnotesize\normalNumber\color{#3}| ulm university  }}]
      (-1.6,0) arc [start angle=-180, end angle=00, radius=1.6];
       \draw[very thick, color=#3] (0,0) circle (1.25);
       \draw[thick, color=#3] (0,0) circle (1.75);
       \draw[thin, color=#3] circle (1.2);
       \foreach \r in {0,20,...,360} {
           \filldraw[#3] (\r:1.2) circle (0.01) node (\r) {};
       }
       \foreach \r in {0,20,...,360} {
           \foreach \rr in {0,20,...,\r}{
               \draw[very thin] (\r) -- (\rr);
           }
       }
       \node[opacity=0.5] at(0.02+#9,-#9) {\fontsize{#6}{16pt}\selectfont#2};
       \draw (0.02,0) node {\fontsize{#6}{16pt}\selectfont\color{#4}#2};
   \end{tikzpicture}
}

\DeclareDocumentCommand{\LILLYxFACULTYxMATHE}{ O{MudWhite} O{FacultyMathexColor} O{} }{%
\LILLYxGENxFACULTY{Mathematik}{$\pi$}{#1}[#2][#3][80pt][\smallNumber][0.125]
}
\DeclareDocumentCommand{\LILLYxFACULTYxPRAKTISCHEINFORMATIK}{ O{MudWhite} O{FacultyPraktischeInformatikxColor} O{} }{%
    \LILLYxGENxFACULTY{Praktische Informatik}{\hspace*{-0.025em}\bfseries\faCode}{#1}[#2][#3][50pt][\smallNumber][0.125]
}
\DeclareDocumentCommand{\LILLYxFACULTYxTHEORETISCHEINFORMATIK}{ O{MudWhite} O{FacultyTheoretischeInformatikxColor} O{} }{%
    \LILLYxGENxFACULTY{Theoretische Informatik}{\bfseries\hspace*{-0.1em}\faCogs}{#1}[#2][#3][45pt][\smallNumber][0.125]
}
\DeclareDocumentCommand{\LILLYxFACULTYxTECHNISCHEINFORMATIK}{ O{MudWhite} O{FacultyTechnischeInformatikxColor} O{} }{%
    \LILLYxGENxFACULTY{Technische Informatik}{\hspace*{-0.025em}\begin{tikzpicture}\filldraw[rounded corners=4pt] (0,0) rectangle ++(1,1);%
        \foreach \y in {1.117,-0.317} {%
            \foreach \x in {0,0.18,...,0.9} {%%
                \filldraw[rounded corners=0.04] (\x,\y) rectangle ++(0.07,0.2); % TOP/BOT
                \filldraw[rounded corners=0.04] (\y,\x) rectangle ++(0.2,0.07); % LEFT/RIGHT
            }%
        }%
    \end{tikzpicture}}{#1}[#2][#3][50pt][\smallNumber][0.125][0.03]
}


\makeatletter


%% Pfad zu den Einstellungen %% Pfad zur Grafik
\DeclareDocumentCommand{\LILLYxTITLE}{ O{\LILLYxPATHxDATA/Semester/\LILLYxSemester/Definitions/GENERAL.tex} O{\LILLYxPATHxDATA/Semester/\LILLYxSemester/Definitions/\LILLYxVorlesung} O{\LILLYxPATHxDATA/Semester/\LILLYxSemester/Graphics/titleimage\LILLYxVorlesung.pdf} } {
    \ifnum\LILLYxSemester>0
        \input{#1}\input{#2}\begingroup
        \clearpage\begin{titlepage}
            \newgeometry{margin = 0pt}
            {\centering\par\vspace*{7em}\par
                {\raggedright\hspace*{\LILLYxLeftsep}\parbox{\linewidth-\LILLYxLeftsep}{\bfseries\usefont{T1}{qzc}{m}{it}\fontsize{40pt}{16pt}\selectfont\FULLTITLE}\par}
                \par\vspace*{11em}\par
                \includegraphics[height=20em]{#3}
                \par\vfill\vspace{2em}\par
                {\raggedright\hspace*{\LILLYxLeftsep}\parbox{\linewidth-\LILLYxLeftsep}{{\scshape \tiny{\LILLY@Typ{} von}\\\large\AUTHOR\\[-0.35\baselineskip]{\normalfont\LILLYxlstTypeWriter\tiny\AUTHORMAIL}\par}
                \vspace{0.5cm}
                {\itshape\small Version vom: \\ \heute\par}}}
            }
            \hbox{}\vfill
            \begin{tikzpicture}
                %\node[above right] at (\LILLYxLeftsep,3.05) {\centering\fontsize{3pt}{2pt}\selectfont\centering Fassung vom \heute{} um \currenttime Uhr\par};
                \draw[thick,\LILLYxFACULTYxCOLOR] (0,3.1) -- ++ (0.6\paperwidth-0.1,0) -- ++(0.1\paperwidth,1) -- ++(0.1\paperwidth+0.1,0) --  ++(0.1\paperwidth,-1) -- ++(0.1\paperwidth,0) ;
                \draw[thin,\LILLYxFACULTYxCOLOR] (0,3.175) -- ++ (0.6\paperwidth-0.175,0) -- ++(0.1\paperwidth,1) -- ++(0.1\paperwidth+0.175,0) --  ++(0.1\paperwidth,-1) -- ++(0.1\paperwidth,0) ;
                \filldraw[\LILLYxFACULTYxCOLOR] (0,0) -- ++(\paperwidth,0) -- ++(0,3) --  ++(-0.1\paperwidth,0) -- ++(-0.1\paperwidth,1) -- ++(-0.1\paperwidth,0) --  ++(-0.1\paperwidth,-1) -- ++(-0.6\paperwidth,0) -- cycle;
                \begin{scope}[every node/.style={color=LightGray}]
                    \node at(0.75\paperwidth,2) {\LILLYxFACULTY};
                    \node[right] at(\LILLYxLeftsep,1.5) {\parbox{0.46\linewidth}{\footnotesize{\LILLYxFlavourText\par
                }}};
                    \node[above left] at(\paperwidth,0.02) {\begingroup\usefont{U}{Konanur}{xl}{n}LILLY\endgroup};
                    \node[above left] at(\paperwidth,0.52) {\tiny \textsf{Ver. \LILLYxVERSION}};
                \end{scope}
            \end{tikzpicture}
        \end{titlepage}\endgroup\restoregeometry
    \else\fi
}

\begin{document}
    \foreach \vl/\se in {ANA1/2,PDP/2,PVS/2,GDBS/2,GDRA/1,FG/1,LAII/1,EIDI/1}{
        \renewcommand{\LILLYxVorlesung}{\vl}
        \renewcommand{\LILLYxSemester}{\se}
        \LILLYxTITLE
    }

    Information! Dieses Dokument wurde generiert durch:
    \begin{lstlisting}[language=lLatex]
\begin{document}
    \foreach \vl/\se in {ANA1/2,PDP/2,PVS/2,GDBS/2,GDRA/1,FG/1,LAII/1,EIDI/1}{
        \renewcommand{\LILLYxVorlesung}{\vl}
        \renewcommand{\LILLYxSemester}{\se}
        \LILLYxTITLE
    }
\end{document}
    \end{lstlisting}

    Der Titlepage-Provider Philosopher lässt sich also dynamisch konfigurieren. Es stehen bisher folgende Siegel zur Verfügung:
    \begin{center}
        \begin{tabular}{*{2}{m{0.2\linewidth}>{\centering\arraybackslash}m{0.25\linewidth}@{\hspace*{2.5em}}}}
            \toprule
                Name & Siegel & Name & Siegel\\
            \midrule
                Mathe & \LILLYxFACULTYxMATHE[black] & Theoretische Inf.
                        & \LILLYxFACULTYxTHEORETISCHEINFORMATIK[black] \\
                Praktische Inf. & \LILLYxFACULTYxPRAKTISCHEINFORMATIK[black] & Technische Inf.
                        & \LILLYxFACULTYxTECHNISCHEINFORMATIK[black] \\
            \bottomrule
        \end{tabular}
    \end{center}
    Geplant sind noch: \T{(Schwerpunkt) Informatik} und \T{Seminar}. Dies sollte alle im Studiengang Informatik enthaltenen Module abdecken.

    Der Code für die Tabelle lautet: %% using ugly !**! for } kw-detection
{\footnotesize\begin{lstlisting}[language=lLatex]
\begin{tabular}{*{2}{m{0.2\linewidth}>{!**!\centering!**!\arraybackslash!**!}m{0.25\linewidth}@{\hspace*{2.5em}}}}
    \toprule
        Name & Siegel & Name & Siegel\\
    \midrule
        Mathe & \LILLYxFACULTYxMATHE[black] & Theoretische Inf.
                & \LILLYxFACULTYxTHEORETISCHEINFORMATIK[black] \\
        Praktische Inf. & \LILLYxFACULTYxPRAKTISCHEINFORMATIK[black] & Technische Inf.
                & \LILLYxFACULTYxTECHNISCHEINFORMATIK[black] \\
    \bottomrule
\end{tabular}
\end{lstlisting}
}

\clearpage

    Die verwendeten Farben gemäß ihrer Definition lauten übrigens:
\begin{lstlisting}[language=lLatex]
%% Mathe
\definecolor{FacultyMathexColor}{RGB}{086,170,028}
%% Praktische Informatik
\definecolor{FacultyPraktischeInformatikxColor}{RGB}{221,101,015}
%% Theoretische Informatik
\definecolor{FacultyTheoretischeInformatikxColor}{RGB}{002,133,221}
%% Technische Informatik
\definecolor{FacultyTechnischeInformatikxColor}{RGB}{177,2,221}
\end{lstlisting}

Für Jake wurde die folgende (Gepard-)Konfiguration verwendet:
\begin{lstlisting}[language=lGepard]
operation	=	file_compile
file		=	philosopher.tex

lilly-show-boxname  = false
lilly-nameprefix    = "De++\
                        mo-" !Angeb :D!

debug		= true  ! Just for fun !
\end{lstlisting}



\end{document}